\chapter{Electrolyte Panel}

\section*{What Is This Test?}
The electrolyte panel is a group of blood tests that measures the main electrolytes in the body: sodium (Na\textsuperscript{+}), potassium (K\textsuperscript{+}), chloride (Cl\textsuperscript{-}), and bicarbonate (total CO\textsubscript{2} / HCO\textsubscript{3}\textsuperscript{-}). When combined with glucose, blood urea nitrogen (BUN), and creatinine, the electrolyte panel constitutes the basic metabolic panel (BMP). The addition of albumin, total protein, calcium, total bilirubin, alkaline phosphatase, AST, and ALT to the BMP creates the comprehensive metabolic panel (CMP). Electrolytes are minerals that carry an electrical charge when dissolved in body fluids. They are essential for maintaining cellular membrane potentials, nerve impulse conduction, muscle contraction (including cardiac muscle), fluid balance, acid-base homeostasis, and enzyme function. The concentrations of sodium, potassium, chloride, and bicarbonate in the blood are precisely regulated by the kidneys, endocrine system (aldosterone, ADH, PTH, vitamin D), and the gastrointestinal tract. Disruptions in any electrolyte can cause life-threatening complications including cardiac arrhythmias, neurological dysfunction, and respiratory failure \cite{seifter2014integration}. The electrolyte panel is one of the most frequently ordered laboratory tests in both inpatient and outpatient settings.

\section*{Alternative Names}
Electrolyte Panel; Serum Electrolytes; Lytes; Na/K/Cl/CO2; Electrolyte Profile; BMP Electrolytes; Chem-4 (sodium, potassium, chloride, bicarbonate); Electrolyte Screen; Basic Metabolic Panel Electrolytes

\section*{Principle}
The electrolyte panel measures four analytes, each by a different principle: (1) Sodium (Na\textsuperscript{+}) is measured by ion-selective electrode (ISE) potentiometry using a glass membrane electrode containing aluminum silicate (NAS glass). The potential difference across the membrane is proportional to the logarithm of sodium activity. (2) Potassium (K\textsuperscript{+}) is measured by ISE using a valinomycin-based membrane. Valinomycin is a cyclic ionophore antibiotic with approximately 10,000:1 selectivity for potassium over sodium. (3) Chloride (Cl\textsuperscript{-}) is measured by ISE with a chloride-sensitive membrane containing quaternary ammonium salts. (4) Total CO\textsubscript{2} (bicarbonate) is measured using an enzymatic method with phosphoenolpyruvate carboxylase (PEPC), which converts bicarbonate and phosphoenolpyruvate to oxaloacetate; oxaloacetate is subsequently reduced to malate by malate dehydrogenase with oxidation of NADH to NAD\textsuperscript{+}, and the decrease in absorbance at 340 nm is proportional to bicarbonate concentration. All four analytes, along with the calculated anion gap [Na -- (Cl + HCO\textsubscript{3})], are measured simultaneously within minutes on a single automated chemistry analyzer platform using integrated ISE modules and spectrophotometric detection. The sample is typically serum (gold-top or red-top tube) or lithium heparin plasma (green-top tube), and is processed via indirect ISE after dilution on most central laboratory analyzers, or via direct ISE on blood gas analyzers and point-of-care devices \cite{tietz2023textbook}.

\section*{History}
The electrolyte panel as a discrete test battery emerged with the development of multi-channel automated chemistry analyzers in the 1960s-1970s. The Technicon SMA 6/60 (Sequential Multiple Analyzer, 6 channels, 60 samples per hour) was introduced in 1965 and included sodium, potassium, chloride, and CO\textsubscript{2} among its measured analytes. This represented a transformative advance from the previous practice of ordering each electrolyte individually. The Chem-4 (later Chem-7, then BMP) became standard in emergency departments, ICUs, and medical-surgical wards, significantly reducing turnaround times and clinical decision-making delays. The basic metabolic panel was codified as a recognized test panel in the 1980s by the American Medical Association's CPT coding system. Point-of-care electrolyte testing was introduced in the 1990s with handheld ISE devices (i-STAT, in 1992, was the first FDA-cleared handheld blood analyzer for electrolytes, glucose, and blood gases). Contemporary chemistry analyzers (Roche Cobas, Abbott Architect, Siemens Atellica) can process over 1,000 electrolyte panels per hour with turnaround times under 10 minutes.

\section*{Why This Test Is Ordered}
\begin{itemize}
  \item Routine health screening as part of the basic metabolic panel (BMP) in annual physical examinations and preoperative assessments
  \item Evaluation of symptoms suggesting electrolyte imbalance: weakness, confusion, seizures, muscle cramps, paresthesias, palpitations, syncope, polyuria, polydipsia, nausea, and vomiting
  \item Monitoring of hospitalized patients receiving IV fluids, total parenteral nutrition (TPN), or medications that significantly affect electrolyte homeostasis (diuretics, ACE inhibitors, ARBs, potassium supplements, corticosteroids, amphotericin B, aminoglycosides, cisplatin, calcineurin inhibitors)
  \item Assessment of renal, cardiac, hepatic, and endocrine disorders that affect fluid and electrolyte balance: chronic kidney disease, congestive heart failure, cirrhosis with ascites, adrenal insufficiency, Conn's syndrome, Cushing's syndrome, diabetes insipidus, SIADH
  \item Evaluation of acid-base disorders: the electrolyte panel provides the bicarbonate value and allows calculation of the anion gap [Na -- (Cl + HCO\textsubscript{3})], which is essential for classifying metabolic acidosis
  \item Preoperative risk stratification in patients with cardiovascular, renal, or hepatic comorbidities
  \item Monitoring of patients on nil-per-os (NPO) status, prolonged nasogastric suction, or large volume diarrhea, where electrolyte losses are significant
  \item Emergency department evaluation of altered mental status, seizures, or cardiac arrest
\end{itemize}

\section*{Primary vs Secondary Test}
The electrolyte panel (with glucose, BUN, and creatinine as the BMP) is a primary, first-line screening test in virtually all clinical settings. It is part of the standard admission order set for hospitalized patients and is often ordered daily in ICU and step-down unit patients. Abnormal results frequently guide initial management decisions and trigger further diagnostic workup.

\section*{How This Test Is Performed}
A healthcare professional will draw blood from a vein, typically in the antecubital fossa (inner elbow). A tourniquet is applied briefly, the skin is cleaned with an alcohol swab, and a needle is inserted. Approximately 3--5 mL of blood is collected into a serum separator tube (gold-top, SST) or a lithium heparin tube (green-top). The tourniquet should be released as soon as blood flow is established---prolonged tourniquet time (>1 minute) can cause hemoconcentration, falsely elevating sodium, chloride, and to a lesser extent potassium. The sample is sent to the laboratory, centrifuged, and the serum or plasma is analyzed on an automated chemistry analyzer. Results for all four electrolytes are typically available within 30--60 minutes. In emergency or ICU settings, a whole blood sample can be run on a point-of-care device (i-STAT, epoc, Gem Premier) or blood gas analyzer using direct ISE, with results available in 2--5 minutes. These point-of-care results correlate well with central laboratory values but may differ slightly because direct ISE measures electrolyte activity in the plasma water phase and is not subject to pseudohyponatremia from hyperlipidemia or hyperproteinemia.

\section*{What Preparation Is Needed}
No fasting is required for the electrolyte panel alone. However, if the panel is ordered as part of a comprehensive metabolic panel (CMP) or lipid profile, 8--12 hour fasting is recommended. Specific patient preparation considerations: (1) Avoid vigorous exercise immediately before the blood draw, as muscle contraction releases potassium and lactic acid (which lowers bicarbonate). (2) Inform the healthcare team of all medications: diuretics (thiazide, loop, potassium-sparing), ACE inhibitors, ARBs, NSAIDs, corticosteroids, insulin, beta-blockers, lithium, desmopressin, SSRIs, and potassium supplements. (3) If receiving IV fluids, the blood should ideally be drawn from the opposite arm to avoid contamination or dilution from the IV infusion. (4) Avoid excessive water intake in the hour preceding the test. (5) Prolonged tourniquet application during phlebotomy should be avoided. (6) For potassium, avoid fist clenching repeatedly during the draw (can increase potassium by 0.5--1.0 mEq/L). (7) Samples must be processed within 4 hours at room temperature or 24 hours refrigerated. For potassium specifically, centrifugation within 30 minutes is ideal to prevent pseudohyperkalemia from cellular potassium leakage.

\section*{Contraindications and Precautions}
No absolute contraindications. Critical pre-analytical factors that can invalidate the electrolyte panel: (1) Hemolysis: the most common cause of pseudohyperkalemia (red blood cell potassium is ~100 mEq/L); hemolyzed specimens should be rejected and redrawn for potassium, though sodium and chloride are minimally affected. (2) EDTA tube (purple-top) contamination: EDTA chelates calcium and magnesium and contains potassium salts. If blood is drawn into an EDTA tube and then poured into a chemistry tube, potassium will be massively elevated (often >10 mEq/L), calcium will be very low, and magnesium may be altered. The standard draw order of tubes must be followed (chemistry tubes before EDTA tubes). (3) Sample drawn from an arm with a running IV: massive dilution, spurious elevation, or contamination depending on the IV fluid composition. (4) Pseudohyponatremia: with indirect ISE methods, severe hypertriglyceridemia (>1500 mg/dL) or hyperproteinemia (>10 g/dL) displaces plasma water, causing apparently low sodium concentration relative to total sample volume but normal sodium in the plasma water phase. Direct ISE (blood gas analyzer, point-of-care) is not affected. (5) Thrombocytosis (>1,000,000/$\mu$L): potassium released from platelets during clotting can elevate serum potassium by 0.5--2.0 mEq/L; plasma lithium heparin samples are preferred in these patients. (6) Leukocytosis (>100,000/$\mu$L): fragile leukemic cells may release potassium into the sample during processing. (7) Delayed centrifugation: prolonged contact between serum/plasma and cells allows potassium leakage from cells and glucose consumption. (8) Cooling before centrifugation: cold temperatures inhibit Na\textsuperscript{+}/K\textsuperscript{+}-ATPase, causing potassium to leak out of cells (do not refrigerate whole blood samples for potassium). (9) Uncapped tubes: CO\textsubscript{2} gas loss to atmosphere lowers measured bicarbonate, causing a spurious low TCO\textsubscript{2} and a falsely elevated calculated anion gap. (10) Bromide or iodide exposure: both ions interfere with the chloride ISE, causing pseudohyperchloremia and a false low or negative anion gap.

\section*{Risks}
Minimal risk from standard venipuncture. Mild pain or bruising at the puncture site is common and resolves within 1--2 days. Rare complications: hematoma, infection (extremely rare), inadvertent arterial puncture (recognized by bright red pulsatile blood; requires prolonged direct pressure), nerve injury (rare, usually transient paresthesia in median or lateral antebrachial cutaneous nerve distribution), and vasovagal syncope (fainting), which occurs in approximately 1--3\% of blood draws. Blood volume drawn (3--5 mL) is negligible even in most anemic patients. In patients with severe thrombocytopenia or coagulopathy, longer direct pressure after venipuncture may be required to prevent hematoma.

\section*{What Happens After the Test}
Pressure is applied with cotton gauze for 1--2 minutes and a bandage is placed. Results for the full electrolyte panel are typically available within 30--60 minutes from the central laboratory, or 2--5 minutes from point-of-care devices. The results of the four electrolytes are interpreted together, not in isolation. The calculated anion gap [Na -- (Cl + HCO\textsubscript{3})] is automatically reported by most laboratory information systems \cite{wallach2022interpretation}. Critical electrolytes: sodium <120 or >160 mEq/L, potassium <2.5 or >6.5 mEq/L, bicarbonate <10 or >40 mEq/L require immediate clinical intervention---these critical values trigger mandatory direct communication from the laboratory to the ordering provider per CLIA and CAP regulations \cite{tierney2023current}. Following an abnormal electrolyte panel, the healthcare provider will: (1) assess volume status, (2) review medications for causative agents, (3) check the ECG for cardiac effects (especially for potassium abnormalities), (4) order additional testing to further characterize the underlying disorder, and (5) initiate specific electrolyte replacement or reduction therapy as indicated. Serial electrolyte monitoring every 4--12 hours is standard in actively managed electrolyte disorders.

\section*{Understanding Results}

\subsection*{Below Normal}
\begin{itemize}
  \item \textbf{Hyponatremia (Na <135 mEq/L):} Classified by volume status---hypovolemic (GI losses, diuretics, adrenal insufficiency, cerebral salt wasting), euvolemic (SIADH, hypothyroidism, psychogenic polydipsia, reset osmostat), or hypervolemic (heart failure, cirrhosis, nephrotic syndrome, CKD). Euvolemic hyponatremia with urine osmolality >100 mOsm/kg and urine sodium >20--40 mEq/L is most commonly SIADH. Hyponatremia is the most common electrolyte abnormality, seen in 15--30\% of hospitalized patients. Na <125 mEq/L: risk of cerebral edema, headache, nausea. <120 mEq/L: seizures, coma \cite{sterns2015disorders}.
  \item \textbf{Hypokalemia (K <3.5 mEq/L):} Renal losses: diuretics (thiazide, loop), primary hyperaldosteronism, Cushing's syndrome, Bartter/Gitelman syndromes, Liddle syndrome, RTA types I and II, amphotericin B, cisplatin, hypomagnesemia (renal potassium-wasting refractory to potassium replacement until magnesium corrected). GI losses: vomiting, NG suction (hypokalemia is primarily from renal potassium loss due to metabolic alkalosis, not from gastric fluid loss which contains only 5--10 mEq/L K), diarrhea (10--80 mEq/L K loss), laxative abuse, villous adenoma. Transcellular shift: insulin, beta-2 agonists, metabolic alkalosis, refeeding syndrome, hypokalemic periodic paralysis. K <3.0 mEq/L: ECG shows flattened T waves, prominent U waves, ST depression, prolonged QT. <2.5 mEq/L: respiratory muscle weakness, rhabdomyolysis, ileus, cardiac arrest risk \cite{gennari1998hypokalemia}.
  \item \textbf{Hypochloremia (Cl <96 mEq/L):} Usually tracks with sodium concentration and is rarely an isolated finding. Metabolic alkalosis with low urine chloride (<10--20 mEq/L) indicates chloride-responsive vomiting/NG suction/diuretic-induced alkalosis. Diuretics, Bartter/Gitelman syndromes, chronic respiratory acidosis (compensatory decrease in chloride as bicarbonate rises).
  \item \textbf{Low bicarbonate (HCO\textsubscript{3} <22 mEq/L):} Metabolic acidosis, classified by anion gap. High-AG (>12 mEq/L): MUDPILES---Methanol, Uremia, DKA, Propylene glycol, Isoniazid/Iron, Lactic acidosis (most common), Ethylene glycol, Salicylates \cite{kraut2007serum}. Normal-AG (hyperchloremic, 8--12 mEq/L): diarrhea, renal tubular acidosis (types I, II, IV), saline infusion, ureteral diversion. Bicarbonate <10 mEq/L: severe acidosis, often life-threatening.
  \item \textbf{Low anion gap (<3 mEq/L):} Uncommon and most often due to hypoalbuminemia (albumin is the major unmeasured anion; for every 1 g/dL decrease in albumin below 4.0, add 2.5 mEq/L to the calculated AG). Other causes: laboratory error, bromide toxicity (pseudohyperchloremia---negative AG), lithium toxicity, IgG multiple myeloma (cationic paraprotein).
\end{itemize}

\subsection*{Above Normal}
\begin{itemize}
  \item \textbf{Hypernatremia (Na >145 mEq/L):} Always reflects hyperosmolality and a relative deficit of water. Hypovolemic: GI losses (diarrhea, vomiting, NG suction, osmotic diarrhea), renal losses (loop diuretics, osmotic diuresis with hyperglycemia or mannitol), excessive sweating. Euvolemic: diabetes insipidus (central from pituitary/hypothalamic pathology; nephrogenic from lithium, hypercalcemia, hypokalemia, genetic mutations). Hypervolemic: iatrogenic (hypertonic saline, sodium bicarbonate infusion), mineralocorticoid excess (hyperaldosteronism, Cushing's), seawater ingestion. Na >155 mEq/L: neurological symptoms prominent. >160 mEq/L: mortality approaches 40--60\%.
  \item \textbf{Hyperkalemia (K >5.0 mEq/L):} First exclude pseudohyperkalemia (hemolysis, prolonged tourniquet, fist clenching, thrombocytosis, leukocytosis, delayed processing). True hyperkalemia: (a) Renal failure (most common---CKD stage 4--5, AKI). (b) Medications: ACE inhibitors, ARBs, spironolactone, eplerenone, amiloride, triamterene, NSAIDs, trimethoprim, tacrolimus, cyclosporine, heparin, succinylcholine. (c) Adrenal insufficiency (Addison's disease, primary hypoaldosteronism). (d) Transcellular shift: metabolic acidosis, insulin deficiency (DKA), rhabdomyolysis, tumor lysis syndrome, massive transfusion. K >6.0 mEq/L: ECG with peaked T waves, prolonged PR. >6.5 mEq/L: widened QRS, sine wave, ventricular fibrillation, asystole---life-threatening emergency.
  \item \textbf{Hyperchloremia (Cl >106 mEq/L):} Hyperchloremic metabolic acidosis (diarrhea, RTA, saline infusion). Primary hyperparathyroidism (mild hyperchloremic acidosis). Dehydration (proportional rise with sodium). Pseudohyperchloremia from bromide/iodide exposure.
  \item \textbf{Elevated bicarbonate (HCO\textsubscript{3} >28 mEq/L):} Metabolic alkalosis. Chloride-responsive (urine Cl <10--20 mEq/L): vomiting, NG suction, diuretic use, respond to volume/NaCl repletion. Chloride-resistant (urine Cl >20 mEq/L): hyperaldosteronism, Cushing's, Bartter/Gitelman, severe hypokalemia, licorice. HCO\textsubscript{3} >40 mEq/L: severe alkalosis, risk of arrhythmia.
  \item \textbf{Elevated anion gap (>12 mEq/L with corrected albumin):} High-AG metabolic acidosis. Gap >20 mEq/L: strongly suggests DKA, toxic alcohol, or severe lactic acidosis. Delta-delta analysis: (measured AG -- 12) + measured HCO\textsubscript{3} should approximate 24--28 mEq/L. If sum >28, concurrent metabolic alkalosis. If sum <22, concurrent normal-AG metabolic acidosis.
\end{itemize}

\section*{Normal Results}
\begin{itemize}
  \item \textbf{Sodium (Na\textsuperscript{+}):} 135--145 mEq/L
  \item \textbf{Potassium (K\textsuperscript{+}):} 3.5--5.0 mEq/L (plasma: 0.1--0.4 mEq/L lower than serum due to platelet release during clotting)
  \item \textbf{Chloride (Cl\textsuperscript{-}):} 96--106 mEq/L
  \item \textbf{Bicarbonate (HCO\textsubscript{3}\textsuperscript{-} / total CO\textsubscript{2}):} 22--28 mEq/L
  \item \textbf{Anion gap:} 8--12 mEq/L [Na -- (Cl + HCO\textsubscript{3})]
  \item \textbf{Albumin-corrected anion gap:} AG + 2.5 $\times$ (4.0 -- measured albumin in g/dL)
  \item \textbf{Critical sodium:} <120 or >160 mEq/L
  \item \textbf{Critical potassium:} <2.5 or >6.5 mEq/L
  \item \textbf{Critical bicarbonate:} <10 or >40 mEq/L \cite{fischbach2023manual}
\end{itemize}

\section*{Follow-up Tests}
\begin{itemize}
  \item \textbf{Comprehensive metabolic panel (CMP)/Basic metabolic panel:} When only electrolytes were drawn, a CMP adds BUN, creatinine, glucose, calcium, albumin, total protein, and liver function tests that are essential to contextualize electrolyte abnormalities and identify underlying renal or hepatic dysfunction.
  \item \textbf{ECG (12-lead):} Urgently if potassium is <3.0 or >6.0 mEq/L. Look for U waves, flattened T waves, prolonged QT (hypokalemia); peaked T waves, prolonged PR, widened QRS, sine wave pattern (hyperkalemia).
  \item \textbf{Arterial blood gas (ABG):} When bicarbonate is abnormal: provides pH, pCO\textsubscript{2}, and calculated HCO\textsubscript{3} for complete acid-base characterization and assessment of respiratory compensation.
  \item \textbf{Plasma osmolality (measured):} If hyponatremia or hypernatremia present. Determine if hyponatremia is hypotonic, isotonic (pseudohyponatremia), or hypertonic (hyperglycemia, mannitol).
  \item \textbf{Urine osmolality and urine sodium:} Essential for SIADH diagnosis (urine osmolality >100 mOsm/kg, urine Na >20--40 mEq/L) and to differentiate causes of hyponatremia by volume status.
  \item \textbf{Serum magnesium:} In refractory hypokalemia---hypomagnesemia causes renal potassium wasting that is refractory to potassium replacement until magnesium is corrected to >2.0 mg/dL.
  \item \textbf{Serum lactate:} Most common cause of elevated anion gap metabolic acidosis (normal <2 mmol/L).
  \item \textbf{Serum ketones (beta-hydroxybutyrate):} Suspected DKA or alcoholic ketoacidosis in high-AG metabolic acidosis.
  \item \textbf{Toxic alcohol screen (methanol, ethylene glycol) and serum osmolal gap:} Unexplained high-AG metabolic acidosis with elevated osmolal gap (>10 mOsm/kg).
  \item \textbf{Plasma renin activity and aldosterone:} For suspected hyperaldosteronism (hypertension, hypokalemia, metabolic alkalosis) or hypoaldosteronism (hyperkalemia, hyponatremia, Type IV RTA).
  \item \textbf{Morning cortisol / ACTH stimulation test:} Suspected adrenal insufficiency causing hyponatremia and hyperkalemia.
\end{itemize}

\section*{References}
\bibliographystyle{unsrtnat}
\bibliography{references}

\clearpage
\section*{Regulatory Status and Authorization Date}
Regulatory status applies to the specific assay, instrument, manufacturer, and intended use rather than to the test name alone. No single approval date can be assigned to this test category. For a commercial assay, verify the current product in the relevant regulator's database: the U.S. Food and Drug Administration (FDA) clearance, approval, or authorization; India's Central Drugs Standard Control Organisation (CDSCO) licence or permission; the European Union In Vitro Diagnostic Regulation (EU IVDR) CE certificate issued through the applicable conformity-assessment route; or the United Kingdom Medicines and Healthcare products Regulatory Agency (MHRA) registration and UKCA/CE status. Laboratory-developed tests may not have an individual FDA, CDSCO, EU IVDR, or MHRA approval date.
\clearpage
